% --- CAPÍTULO 1 ---
\chapter{Introdução}
O diabetes mellitus é uma das doenças crônicas mais prevalentes no mundo, afetando milhões de pessoas e representando uma preocupação crescente para os sistemas de saúde pública \cite{SBD2025, WHO2016, IDF2021}. Seu diagnóstico precoce e tratamento adequado são essenciais para evitar complicações graves, como doenças cardiovasculares, falência renal, retinopatia e neuropatia. Dentre os tipos existentes, o diabetes tipo 2 é o mais comum, estando frequentemente associado a fatores comportamentais, como má alimentação, sedentarismo e obesidade \cite{IDF2021}.

Com o avanço da tecnologia e o crescimento exponencial da disponibilidade de dados médicos e comportamentais, surgem novas oportunidades para apoiar o diagnóstico e a prevenção de doenças por meio da inteligência artificial. Entre as abordagens mais promissoras nesse campo estão as redes neurais artificiais, que integram o campo do aprendizado de máquina e se destacam por sua capacidade de identificar padrões complexos em grandes conjuntos de dados \cite{Miotto2018,Deo2015}.

Esses modelos podem analisar variáveis como idade, índice de massa corporal (IMC), níveis de glicose, histórico familiar e hábitos de vida, permitindo a criação de sistemas capazes de estimar o risco de um indivíduo desenvolver diabetes com alta precisão \cite{Sisodia2018,Kavakiotis2017}. Essa capacidade torna as redes neurais uma ferramenta valiosa no apoio à decisão clínica, possibilitando diagnósticos mais rápidos, assertivos e estratégias de intervenção precoce.

É neste contexto que se insere o presente trabalho, ao propor a aplicação de redes neurais artificiais como uma solução tecnológica para a predição de casos de diabetes, contribuindo para a detecção precoce e melhoria na qualidade de vida dos pacientes.

\section{Problematização}
O diabetes tipo 2 representa atualmente um dos maiores desafios de saúde pública em todo o mundo, devido ao seu crescimento acelerado, alto potencial de morbidade e aos custos associados ao tratamento de complicações crônicas. \cite{SBD2025,WHO2016,IDF2021}. Apesar do conhecimento dos fatores de risco e dos avanços nos métodos de detecção, ainda há dificuldade significativa na identificação precoce de indivíduos propensos ao desenvolvimento da doença, especialmente em contextos de grande demanda e recursos limitados. \cite{IDF2021} Essa limitação reforça a necessidade de buscar soluções inovadoras que auxiliem na triagem, no diagnóstico precoce e na prevenção do diabetes tipo 2, permitindo melhorar a qualidade de vida dos pacientes e otimizar os recursos do sistema de saúde. \cite{Sisodia2018,Kavakiotis2017}

Nesse cenário, a integração de tecnologias baseadas em inteligência artificial se apresenta como uma alternativa promissora para superar desafios atuais e apoiar profissionais da saúde na tomada de decisão clínica.

\section{Solução Proposta}
Diante dessa realidade, este trabalho propõe o desenvolvimento e a avaliação de um modelo preditivo baseado em redes neurais artificiais. A intenção é utilizar dados clínicos e comportamentais acessíveis para produzir uma solução computacional capaz de auxiliar profissionais de saúde na identificação precoce de indivíduos sob risco de desenvolver diabetes tipo 2. O trabalho visa comparar modelos que utilizam técnicas modernas de aprendizado de máquina com algoritmos clássicos, buscando aumentar a precisão, agilidade e acessibilidade do apoio ao diagnóstico.

\subsection{Delimitação do Escopo}
Este estudo foca exclusivamente na aplicação de redes neurais para a predição de diabetes tipo 2 em adultos, utilizando um conjunto de dados secundários públicos, com variáveis essencialmente clínicas e comportamentais. O trabalho não se propõe a substituir o diagnóstico médico, mas sim oferecer uma ferramenta complementar para triagem e apoio à decisão clínica, respeitando os limites éticos e técnicos do uso da inteligência artificial na saúde.

\subsection{Justificativa}
A relevância do tema se justifica tanto pelo impacto epidemiológico do diabetes quanto pelo potencial das tecnologias de inteligência artificial em aprimorar processos de prevenção e diagnóstico em saúde pública. O desenvolvimento de modelos acessíveis e escaláveis pode contribuir para otimizar recursos, ampliar o alcance da triagem preventiva e melhorar o prognóstico dos pacientes devido à intervenção precoce. O estudo também proporciona aprendizado aplicado, integrando ciência de dados, programação e fundamentos de saúde.

\section{Objetivos}
Este trabalho visa desenvolver e avaliar um modelo de rede neural artificial capaz de predizer casos de diabetes tipo 2 a partir da análise de dados clínicos e comportamentais, com o intuito de apoiar o diagnóstico precoce e contribuir para estratégias de prevenção.

\subsection{Objetivo Geral}
Desenvolver e avaliar um modelo de rede neural artificial capaz de predizer casos de diabetes tipo 2 com base em dados clínicos e comportamentais, de modo a apoiar a triagem e o diagnóstico precoce em ambientes de saúde.

\subsection{Objetivos Específicos}
Este trabalho possui como objetivos específicos:
\begin{itemize}
    \item Realizar levantamento e análise do estado da arte sobre o uso de redes neurais e modelos de machine learning para predição de doenças crônicas.
    \item Pré-processar e analisar o conjunto de dados utilizado, identificando os atributos mais relevantes ao risco do diabetes tipo 2.
    \item Implementar e treinar modelos de rede neural artificial, ajustando parâmetros e avaliando desempenho.
    \item Comparar a precisão do modelo proposto com métricas clássicas e métodos alternativos, discutindo limitações e potencial de aplicação.
    \item Sugerir recomendações para aprimoramento da ferramenta e sua eventual integração em fluxos clínicos reais.
\end{itemize}